\section{A/B updates}

\begin{frame}{A/B updates}
  \begin{itemize}
    \item Mecanism for Over-The-Air (OTA) \textit{image} updates
    \item Key idea: avoid one update bricking the device
    \begin{itemize}
      \item Even if tested, one update might prevent the system from booting
      \item Enable the device to recover from a bad update automatically
      \item Avoid RMAs or in-the-field intervention
    \end{itemize}
    \item The key idea is to maintain 2 copies of the system, in two \textbf{slots}
    \begin{itemize}
      \item minimum is 2 rootFS
      \item usually also includes the kernel
      \item can involve the bootloader, depending on the bootROM
    \end{itemize}
    \item Non-A/B OTAs are \href{https://source.android.com/docs/core/ota/nonab}{deprecated in Android}
          since Android 6 (pre August 2016)
  \end{itemize}
\end{frame}

\begin{frame}{A/B updates: slots}
  \begin{itemize}
    \item A/B systems have 2 slots: slot A and slot B
    \item They each have dedicated boot locations (usually boot medium partitions)
    % Because OTA happen on live systems, there is always a current slot
    \item On the running system, the slot booted is the \textbf{current} slot
    \item An OTA should never touch boot locations from the \textbf{current} slot
    \item Each slot can be marked \textbf{bootable}
    \item Updates will usually be built independently from the slot they will occupy
  \end{itemize}
\end{frame}

\begin{frame}{A/B updates: strategy}
  \begin{itemize}
    \item At boot, the system detect which slot (A or B) is the \textbf{current} slot
    \item The updater will locate the boot locations for the \textbf{alternate} slot
    \item The updater then applies the update to the \textbf{alternate} boot locations
    \item The updater marks the \textbf{alternate} slot as \textbf{next} to be booted
    \item Eventually, the system reboots. Possibly triggered at the end of the update
    \item At boot, the bootloader detects the \textbf{next} slot and boots it
    \begin{itemize}
      \item If the system fails to boot, the watchdog reboots the system. If the slot
            fails to boot several times, the system boots the other slot.
      \item If the boot succeed, the system runs checks and if they pass, marks the
            \textbf{current} slot as \textbf{primary}
    \end{itemize}
  \end{itemize}
\end{frame}

\begin{frame}{A/B updates: bootloader}
  \begin{itemize}
    \item The bootloader is responsible for:
    \begin{itemize}
      \item Booting the \textbf{primary} slot by default
      \item Keeping track of boot failures per slot
      \item Trying out the \textbf{alternate} slot if it is marked as \textbf{next}
      \item Optionally, telling the kernel which slot has been booted (A or B)
    \end{itemize}
    \item Communication between bootloader and the system is short
  \end{itemize}
\end{frame}


\begin{frame}{Updater: SWUpdate}
  \begin{itemize}
    \item Open source project
    \item \href{https://swupdate.org/}{website} -- \href{https://github.com/sbabic/swupdate}{code} --  \href{https://sbabic.github.io/swupdate/index.html}{doc}
    \item Used e.g. by the \href{https://cip-project.org/}{Civil Infrastructure Project}
    \item Integrated in both Buildroot and Yocto
    \item Written in C
    \item Uses \code{libconfig} syntax for configuration
  \end{itemize}
\end{frame}

\begin{frame}[fragile]
  \frametitle{SWUpdate: configuration}
  \begin{minted}[fontsize=\tiny]{yaml}
software =
{
        version = "0.1.0";
        description = "Firmware update for XXXXX Project";

        hardware-compatibility: [ "1.0", "1.2", "1.3"];

        partitions: (
                { name = "rootfs"; device = "mtd4"; size = 104896512; },
                { name = "data"; device = "mtd5"; size = 50448384; }
        );

        images: (
                { filename = "rootfs.ubifs"; volume = "rootfs"; },
                { filename = "swupdate.ext3.gz.u-boot"; volume = "fs_recovery"; },
                { filename = "sdcard.ext3.gz"; device = "/dev/mmcblk0p1"; compressed = "zlib";},
                { filename = "bootlogo.bmp"; volume = "splash"; },
                { filename = "uImage.bin"; volume = "kernel"; },
                { filename = "fpga.txt"; type = "fpga"; },
                { filename = "bootloader-env"; type = "bootloader"; }
        );

        files: ({ filename = "README"; path = "/README"; device = "/dev/mmcblk0p1"; filesystem = "vfat"; });

        scripts: ( { filename = "erase_at_end"; type = "lua"; }, { filename = "display_info"; type = "lua"; });

        bootenv: (
                { name = "vram"; value = "4M"; },
                { name = "addfb"; value = "setenv bootargs \${bootargs} omapfb.vram=1:2M,2:2M,3:2M omapdss.def_disp=lcd"; }
        );
}
  \end{minted}
\end{frame}

\begin{frame}{Updater: RAUC}
  \begin{itemize}
    \item Robust Auto-Update Controller
    \item Also open source
    \item \href{https://rauc.io/}{website} -- \href{https://github.com/rauc/rauc}{code} -- \href{https://rauc.readthedocs.io/en/latest/}{doc}
    \item Integrated in Buildroot and Yocto
    \item Written in python
    \item Uses a configuration and \code{manifest} in INI format
    \item Updates are handled in the form of \code{bundles}
  \end{itemize}
\end{frame}

\begin{frame}[fragile]
  \frametitle{RAUC: configuration}
  \begin{minted}[fontsize=\scriptsize]{yaml}
[system]
compatible=rauc-demo-x86
bootloader=grub
mountprefix=/mnt/rauc
bundle-formats=-plain

[keyring]
path=demo.cert.pem

[slot.rootfs.0]
device=/dev/sda2
type=ext4
bootname=A

[slot.rootfs.1]
device=/dev/sda3
type=ext4
bootname=B
  \end{minted}
\end{frame}

\begin{frame}[fragile]
  \frametitle{RAUC: Bundle \href{https://rauc.readthedocs.io/en/latest/reference.html\#sec-ref-manifest}{manifest}}
  \begin{itemize}
    \item Must be called \code{manifest.raucm}
  \end{itemize}
  \begin{minted}[fontsize=\scriptsize]{yaml}
[update]
compatible=Test Platform
version=2023.11.0

[bundle]
format=verity

[image.rootfs]
filename=system-image.ext4

[image.bootloader]
filename=barebox.img
  \end{minted}
\end{frame}

\begin{frame}{Updates backend: \href{https://hawkbit.eclipse.dev}{hawkBit}}
  \begin{itemize}
    \item Open Source project from the Eclipse foundation
    \item Back-end for update rollout
    \item Can break down shipping updates into groups
    \item Includes a management interface and visualisation
    \item Supported by both RAUC and SWUpdate
  \end{itemize}
\end{frame}

\subsection{Examples}

\begin{frame}{A/B updates using RAUC + U-Boot}
  \begin{itemize}
    \item U-Boot is natively supported by RAUC
    \item RAUC will use U-Boot's environment to affect boot order
    \begin{itemize}
      \item \code{BOOT_ORDER}: Names of all slots, in order of priority
      \item \code{BOOT_<slot_name>_LEFT}: remaining boot attempts for the slot
    \end{itemize}
    \item The environments variables will be read and set by a
          boot script
    \begin{itemize}
	    \item RAUC includes an \href{https://github.com/rauc/rauc/blob/master/contrib/uboot.sh}{example}
            (see next slide)
    \end{itemize}
    \item Alternatively, one can configure U-Boot to use the
          \href{https://docs.u-boot.org/en/latest/develop/bootstd/rauc.html}{RAUC bootmeth}
  \end{itemize}
\end{frame}

\begin{frame}[fragile]
  \frametitle{A/B updates using RAUC + U-Boot}
  \begin{minted}[fontsize=\tiny]{shell}
test -n "${BOOT_ORDER}" || setenv BOOT_ORDER "A B"
test -n "${BOOT_A_LEFT}" || setenv BOOT_A_LEFT 3
test -n "${BOOT_B_LEFT}" || setenv BOOT_B_LEFT 3

setenv bootargs
for BOOT_SLOT in "${BOOT_ORDER}"; do
  if test "x${bootargs}" != "x"; then
    # skip remaining slots
  elif test "x${BOOT_SLOT}" = "xA"; then
    if test 0x${BOOT_A_LEFT} -gt 0; then
      echo "Found valid slot A, ${BOOT_A_LEFT} attempts remaining"
      setexpr BOOT_A_LEFT ${BOOT_A_LEFT} - 1
      setenv load_kernel "nand read ${kernel_loadaddr} ${kernel_a_nandoffset} ${kernel_size}"
      setenv bootargs "${default_bootargs} root=/dev/mmcblk0p1 rauc.slot=A"
    fi
  elif test "x${BOOT_SLOT}" = "xB"; then
    if test 0x${BOOT_B_LEFT} -gt 0; then
      setexpr BOOT_B_LEFT ${BOOT_B_LEFT} - 1
      setenv load_kernel "nand read ${kernel_loadaddr} ${kernel_b_nandoffset} ${kernel_size}"
      setenv bootargs "${default_bootargs} root=/dev/mmcblk0p2 rauc.slot=B"
    fi
  fi
done

if test -n "${bootargs}"; then
  saveenv
else
  echo "No valid slot found, resetting tries to 3"
  setenv BOOT_A_LEFT 3
  setenv BOOT_B_LEFT 3
  saveenv
  reset
fi

run load_kernel
bootm ${loadaddr_kernel}
  \end{minted}
\end{frame}


\begin{frame}{A/B updates on RPi 5}
  \begin{itemize}
    \item Raspberry Pis have an interesting \href{https://www.raspberrypi.com/documentation/computers/raspberry-pi.html\#boot-sequence}{boot sequence}
    \item the boot actually stats on the VideoCore, which is the GPU
    \item this lets users configure the CPU's bootloader via \code{.txt} files
    \begin{itemize}
      \item \href{https://www.raspberrypi.com/documentation/computers/config_txt.html\#what-is-config-txt}{\code{config.txt}}
      \item \href{https://www.raspberrypi.com/documentation/computers/config_txt.html\#autoboot-txt}{\code{autoboot.txt}}
    \end{itemize}
    \item the RPi bootloader already supports A/B updates, via its 
          \href{https://www.raspberrypi.com/documentation/computers/config_txt.html\#the-tryboot-filter}{\code{tryboot}} feature
    \item It also supports dynamically changing configuration based on the
          \href{https://www.raspberrypi.com/documentation/computers/config_txt.html\#boot_partition-2}{\code{boot_partition}}
    \begin{itemize}
      \item This helps making the OTA update slot-agnostic
    \end{itemize}
  \end{itemize}
\end{frame}

